% Supporting Information body (from COMPASS_JCTC_SUPPORTING_INFORMATION.tex).
% Included in ``paper/COMPASS_BTP_REPORT.tex`` after \appendix; table/figure numbering S-prefix set in main file.

\chapter{Molecular dynamics protocol (source and target trajectories)}

Production trajectories for the wild-type alanine-rich peptide A10 (Ace-(Ala)$_{10}$-NMe) and the mutant A8Q2 (Ace-(Ala)$_8$(Gln)$_2$-NMe) were generated with \textsc{OpenMM}~8.x.\cite{eastman2024openmm8}
We used the additive CHARMM36m protein force field with CHARMM-modified TIP3P water as distributed with \textsc{OpenMM} (\path{charmm36_2024.xml} and matching water parameters).\cite{huang2017charmm36m,kim2015charmmgui}
Nonbonded interactions used a nominal 1.2~nm cutoff with smoothing from 1.0~nm, particle-mesh Ewald electrostatics, constraints on bonds to hydrogen, and rigid water.
Integration used a Langevin middle integrator with 2~fs timestep and friction 1~ps$^{-1}$.
Each system was energy-minimized, equilibrated (NVT with peptide heavy-atom restraints followed by NPT), and simulated in the NVT ensemble for production at \textbf{330~K} to accelerate helix--coil sampling relative to room temperature while preserving the qualitative physics relevant to this benchmark.
The same protocol and temperature were used for A10 and A8Q2 so relative comparisons are internally consistent.

Peptides were solvated in a periodic cubic box with roughly 1.2--1.5~nm buffer and 0.15~M NaCl (see project \path{docs/EXPERIMENT_PROTOCOL.md} for box and ion counts).
Production coordinates were saved at \textbf{0.1~ns} (100~ps) spacing for downstream featurization, consistent with the chunk library described in the main text.
The full A8Q2 trajectory was partitioned into non-overlapping chunks for COMPASS replay; reference observables $f_{\mathrm{helix}}^{\mathrm{ref}}$ and $f_{\mathrm{open}}^{\mathrm{ref}}$ were computed from the same full trajectory.

\chapter{TICA and SPIB training details}

TICA was fit on backbone $\phi/\psi$ (and related) torsional features from the A10 source trajectory using PyEMMA's symmetric time-lagged covariance estimator at a lag matched to the implementation used for production analysis (order 1--4~ns in frame units; see repository \path{src/} and phase-5 analysis logs).
SPIB encoders were trained in PyTorch on the same A10 feature pipeline as in the main experiments, with latent dimension one or two and information-bottleneck trade-off as reported in the main text.

\chapter{Benchmark design and output structure}

The archive-replay benchmark uses pre-recorded A8Q2 trajectory chunks as a library; no new MD is
run during benchmarking.
Chunk scoring and policy dispatch are implemented in \path{src/compass/v2/bench.py}.
The experiment orchestration script \path{scripts/run_experiment_suite.py} resolves caching
dependencies (P1/P2/P3 are CV-agnostic; P2/P3/P4 are noise-agnostic) and dispatches runs via
\texttt{multiprocessing.Pool}.
Reference quantities ($\pi_{\mathrm{ref}}$, per-bin helix fraction, helix--coil boundary TIC-1
value) are precomputed once from the full A8Q2 trajectory by
\path{scripts/precompute_reference.py} and stored under \path{outputs/v2/reference/}.

Output paths follow the convention:
\begin{verbatim}
outputs/v2/experiments/p_{policy}/cv_{cv}/chunk_{c}ns/noise_{n}/repeat_{rr}/
\end{verbatim}
Each \path{v2\_metrics.csv} records per-step: cumulative budget (ns), reweighted helix-fraction
error, reweighted open-fraction error, JSD of the 2D free-energy surface, target-space coverage,
source-space coverage, and cumulative helix$\leftrightarrow$coil transitions.

\chapter{Extended theory, schedule, and supplementary figures}

\subsection*{TICA and SPIB constructions}

Time-lagged independent component analysis seeks coordinates $\boldsymbol\chi(\mathbf{x})$ that maximize autocorrelation at lag $\tau$.\cite{perezhernandez2013tica,schwantes2013improvements}
For mean-free featurized configurations $\mathbf{x}_t$, the instantaneous and symmetric time-lagged covariance matrices are
\begin{equation}
\begin{split}
    \mathbf{C}(0) &= \mathbb{E}_t\!\left[\mathbf{x}_t \mathbf{x}_t^{\top}\right],\\
    \mathbf{C}(\tau) &= \tfrac{1}{2}\mathbb{E}_t\!\left[\mathbf{x}_t \mathbf{x}_{t+\tau}^{\top} + \mathbf{x}_{t+\tau} \mathbf{x}_t^{\top}\right].
\end{split}
    \label{eq:si_tica_covariances}
\end{equation}
\addcontentsline{loe}{eqloe}{\protect\numberline{\theequation}TICA instantaneous and time-lagged covariances}
These define $\mathbf{C}(\tau)\mathbf{u}_i = \lambda_i \mathbf{C}(0)\mathbf{u}_i$, whose top eigenvectors yield the TICA basis used for $p_s(b)$.
The SPIB objective used in training is
\begin{equation}
\begin{split}
    \mathcal{L}_{\mathrm{SPIB}}(\phi,\theta) &=
    \mathbb{E}_{q_\phi}\!\left[\log p_\theta(y_{t+\tau}\mid \mathbf{z}_t)\right]
     - \beta\,D_{\mathrm{KL}}\!\left(q_\phi(\mathbf{z}\mid\mathbf{x})\,\Vert\,r(\mathbf{z})\right).
\end{split}
    \label{eq:si_spib_loss}
\end{equation}
\addcontentsline{loe}{eqloe}{\protect\numberline{\theequation}SPIB variational training objective}

\subsection*{Two-phase component weights}

\begin{table}[htbp]
    \centering
    \caption{Weights $w_k(t)$ for the \texttt{two\_phase} scheduler used in code (main text): three successive segments---exploratory, balanced, and exploit---occupy 30\%, 40\%, and 30\% of the budget, respectively. Redundancy is negative (latent proximity penalty); a bin-fatigue prefactor $1/(1+n_b)$ scales the full combination in Eq.~(\ref{eq:score}).}
    \label{tab:si_policy_weights}
    \small
    \setlength{\tabcolsep}{3.5pt}
    \begin{tabular}{@{}lrrr@{}}
        \toprule
        Component $z_k$ & Expl. & Bal. & Exploit \\
        \midrule
        Source deficit ($d^{IG}_t$) & 0.15 & 1.25 & 1.45 \\
        Source representative & 0.10 & 0.65 & 0.85 \\
        Source-boundary trust & 1.20 & 0.25 & 0.15 \\
        Frontier disagreement & 0.90 & 0.35 & 0.10 \\
        Unvisited & 0.75 & 0.25 & 0.15 \\
        Novelty & 0.65 & 0.25 & 0.10 \\
        Redundancy & $-0.90$ & $-0.80$ & $-0.60$ \\
        \bottomrule
    \end{tabular}
\end{table}

